(1) 在外部理想入流区 \(|E_z|\sim v_{\rm in}B_0\)；在薄电阻片中心对流项较小，\(J_z\sim B_0/(\mu_0\delta)\)，故 \(|E_z|\sim \mu_0\eta_mJ_z\sim\eta_mB_0/\delta\)。因稳态Faraday定律使面外电场近似空间常数，
\[
v_{\rm in}\sim\frac{\eta_m}{\delta}.
\]
不可压的单位深度两侧入流流率与两端出流流率量级匹配：\(2L v_{\rm in}\sim2\delta v_{\rm out}\)，即 \(v_{\rm in}L\sim v_{\rm out}\delta\)。
(2) 上游磁压标度 \(B_0^2/(2\mu_0)\) 释放驱动出流动压 \(\rho v_{\rm out}^2/2\)，取 \(v_{\rm out}\sim v_A=B_0/\sqrt{\mu_0\rho}\)。消去入流速度：
\[
\delta^2\sim\frac{\eta_m L}{v_A},\qquad
\frac{\delta}{L}\sim S^{-1/2},\qquad
\frac{v_{\rm in}}{v_A}\sim S^{-1/2}.
\]
因此 \(|E_z|\sim B_0v_A/\sqrt S\)。所有关系是主导阶标度而非声称精确的单位系数。
(3) 入流动压与磁压之比 \(\rho v_{\rm in}^2/(B_0^2/\mu_0)\sim S^{-1}\)，片纵横比 \(\delta/L\sim S^{-1/2}\)，重联率 \(v_{\rm in}/v_A\sim S^{-1/2}\)。因为 \([\eta_m]=L^2/t\)，\(S\) 无量纲，\([\delta]=L\)，\([v_{\rm in}]=L/t\)，\([E]=[vB]\) 正确。固定 \(L,v_A\) 下 \(\eta_m\to0^+\) 导致层宽与重联率均趋零，这说明层流Sweet--Parker机制在理想MHD极限并不提供快速重联。
(4) 两侧各长约 \(2L\) 还是 \(L\) 取决于半长约定；固定本题每侧总受流长度取 \(2L\)，故两侧单位深度磁Poynting输入率
\[
P_{\rm in}\sim4L\frac{B_0^2}{\mu_0}v_{\rm in}.
\]
两端每端宽约 \(2\delta\)，动能输出率
\(P_{\rm kin}\sim2(2\delta)\frac{\rho v_A^3}{2}=2\delta\rho v_A^3\sim2Lv_{\rm in}B_0^2/\mu_0\)，与输入为同一数量级。磁张力做功、压强焓、喷流剩余磁能以及电阻热的分配依赖真实剖面，不能把 \(J_z=B_0/(\mu_0\delta)\) 当作整片严格常数而求精确比例。超高 \(S\) 可发生撕裂/等离子体团不稳定，并可能进入Hall、动理学或碰撞无关重联区，不应将上述标度外推为通用定律。
